\documentclass[11pt]{article}
\usepackage{amsmath,amssymb,amsthm}
\usepackage[margin=1in]{geometry}

\newcommand{\RespLC}{\mathsf{Resp}_{\mathrm{lc}}}
\newcommand{\SelResp}{\mathsf{Sel}_{\RespLC}}
\newcommand{\Msrc}{\mathcal{M}_{\mathrm{src}}}
\newcommand{\geff}{g_{\mathrm{eff}}}
\newcommand{\XiR}{\Xi_X^R}
\newcommand{\SX}{S_X}
\newcommand{\Splus}{S_X^+}

\newtheorem{proposition}{Proposition}

\title{Resp\_lc Source-Extension Candidate Smuggling Audit}
\author{The AEther-Flow Research Project}
\date{2026-06-18}

\begin{document}
\maketitle

\section*{Control Status}

This artifact is a draft/control Smuggling Auditor packet for
\texttt{RT-20260614-058}. It audits the \(\RespLC\) source-extension
candidate formalized in \texttt{RT-20260614-057}. It is not canonical
ontology, not source-extension adoption, not a canonical \(\RespLC\)
witness, not \(\Msrc\), not \(\geff\), not matter coupling, not
Einstein equations, not benchmark evidence, not a Gate Chair verdict,
and not completed-derivation evidence.

\section{Audit Target}

The audited packet introduced the draft/control response-selector
extension datum
\[
  \XiR=(O_X^R,N_X^R,K_X^R,o_X^R,\nu_X^R,\kappa_X^R),
\]
the expanded source tuple
\[
  \Splus=(\SX,\XiR),
\]
and the candidate selector
\[
  \SelResp^{+}(\Splus)=(o_X^R,\nu_X^R,\kappa_X^R).
\]
It also defined the forgetful reduct \(F_X(\Splus)=\SX\). The present
audit asks whether those definitions hide a target metric, target
proper time, target observer protocol, target calibration, matter
coupling, benchmark evidence, validation authority, Gate Chair
authority, or source-extension adoption.

\section{Audit Screens}

\paragraph{Target metric import.}
The formal dependency list of \(\SelResp^{+}\) is the distinguished
triple inside \(\XiR\). The definition does not call a Lorentzian
metric, manifold metric law, geodesic structure, causal cone, curvature
object, Einstein-equation object, or benchmark solution. This is a
conditional source-purity pass only under the reading that \(O_X^R\),
\(N_X^R\), \(K_X^R\), \(o_X^R\), \(\nu_X^R\), and \(\kappa_X^R\) are
source-language extension symbols.

\paragraph{Proper-time and calibration import.}
The component \(\nu_X^R\) is named as a positive response-normalization
datum. The audit detects no explicit target proper-time scale, physical
ruler calibration, metric normalization, or empirical unit in its
definition. The pass is conditional: if \(\nu_X^R\) were later
identified with target proper time or empirical calibration, the
candidate would fail this screen.

\paragraph{Observer-protocol import.}
The component \(o_X^R\) is introduced as a source-side
response-orientation datum rather than as target time orientation,
target causal orientation, or observer four-velocity data. No target
observer protocol is present in the written dependency graph. The datum
remains an added extension symbol, not a value derived from \(\SX\).

\paragraph{Matter-coupling and token-semantics import.}
The component \(\kappa_X^R\) maps source readout-token classes into
\(K_X^R\). As written, it does not define stress-energy response,
universal matter coupling, detector semantics, or empirical readout
calibration. It may be retained only as uninterpreted source-side token
semantics until further stress testing and any adoption gate.

\paragraph{Benchmark and validation-authority import.}
Registry rows, role names, validation status, handoff labels, and task
continuity are process-control facts. They do not become physics
evidence for \(\RespLC\), \(\Msrc\), \(\geff\), matter coupling,
Einstein equations, benchmark promotion, Gate Chair review, or a
completed derivation.

\paragraph{Source-extension adoption leakage.}
The candidate is a conservative definitional extension candidate over
an expanded source signature. That classification does not imply that
\(\SX\) determines \(\XiR\), that \(\XiR\) is canonical ontology, or
that the extension is admissible for benchmark derivation. Adoption
would require the later audit, Refuter, and protected human-gated
authority sequence.

\section{Dependency Lemmas}

\begin{proposition}[Conditional no-target-import screen]
Assume \(\XiR\) is interpreted only as source-language extension data
and no interpretation map from \(\XiR\) to target metric, target proper
time, target observer protocol, matter coupling, benchmark evidence, or
validation authority is supplied. Then the definition
\[
  \SelResp^{+}(\Splus)=(o_X^R,\nu_X^R,\kappa_X^R)
\]
does not import those target structures by definition.
\end{proposition}

\begin{proof}
The right-hand side of \(\SelResp^{+}\) names only the three
distinguished components of \(\XiR\). The tuple \(\XiR\) contains
source-language symbols and does not mention a target metric tensor,
proper-time functional, observer protocol, calibration rule,
stress-energy tensor, benchmark solution, validation status, role
authority, or Gate Chair verdict. Under the stated assumption, no
definition in the dependency chain can inspect those target structures.
Therefore the written candidate passes this no-target-import screen.
\end{proof}

\begin{proposition}[Audit pass does not imply adoption]
The conditional no-target-import pass for \(\XiR\) does not derive
\(\XiR\) from \(\SX\), does not adopt \(\XiR\) as canonical source
ontology, and does not establish canonical \(\RespLC\).
\end{proposition}

\begin{proof}
The construction explicitly moves from \(\SX\) to the expanded tuple
\(\Splus=(\SX,\XiR)\). The forgetful reduct \(F_X(\Splus)=\SX\)
supports old-language conservativity, but it also shows that the new
components are deleted when the extension is forgotten. Thus the audit
can certify only that the extension is not a hidden target import under
the source-symbol reading. It does not supply a source-only selection
rule for \(\XiR\) from \(\SX\), and it cannot replace the required
adoption gates.
\end{proof}

\section{Audit Verdict}

The audit verdict is:
\[
\texttt{conditional\_source\_purity\_pass\_with\_source\_extension\_adoption\_block}.
\]

No explicit target metric, target proper-time scale, target observer
protocol, target calibration, matter-coupling law, benchmark evidence,
validation-authority import, or Gate Chair authority import is detected
in the written candidate under a source-symbol reading. The pass is
conditional and non-promotional. The candidate remains draft/control
source-extension data. It is not a derivation from the current retained
tuple and is not adopted ontology.

\section{Distance-to-GR Status}

\begin{center}
\begin{tabular}{p{0.33\linewidth}p{0.57\linewidth}}
\hline
Burden & Status \\
\hline
Source ontology primitives & Unchanged; no canonical ontology edit is made and \(\XiR\) remains draft/control extension data. \\
Source equivalence EqSrc & Unchanged; general EqSrc remains undischarged. \\
RetainH & Unchanged; no canonical RetainH adoption is supplied. \\
GenH & Unchanged; no canonical GenH adoption is supplied. \\
ObsLoc\_lc & Unchanged; local exact-branch witness status is not promoted. \\
Resp\_lc & Hidden-import audit passes conditionally for the source-extension candidate; canonical \(\RespLC\) remains blocked pending Refuter stress and any adoption gate. \\
M\_src & Not started; no source manifold construction is supplied. \\
g\_eff & Not started; no effective metric is supplied. \\
Matter coupling & Not started; \(\kappa_X^R\) is not a coupling law. \\
Einstein equations & Not started. \\
Finite-variation robustness & Unchanged; no new finite-variation robustness theorem is supplied. \\
Benchmark promotion & Blocked. \\
Gate Chair status & Blocked and human-gated. \\
Current route freeze or hard-fail status & Source-extension candidate route is not frozen; the audited candidate now requires Refuter stress. \\
\hline
\end{tabular}
\end{center}

\section{Parent-Child Synthesis}

This artifact is the fused parent output for
\texttt{role\_decomposition.mode = parent\_child\_parallel\_synthesis}.
The child outputs are:
\begin{itemize}
  \item \texttt{child\_phys\_math\_resp\_lc\_source\_extension\_candidate\_smuggling\_audit.yaml};
  \item \texttt{child\_phys\_phil\_resp\_lc\_source\_extension\_candidate\_smuggling\_audit.yaml}.
\end{itemize}
The conflict review is
\texttt{parent\_conflict\_review\_resp\_lc\_source\_extension\_candidate\_smuggling\_audit.yaml}
and found no blocking conflict. The fusion notes are
\texttt{parent\_fusion\_notes\_resp\_lc\_source\_extension\_candidate\_smuggling\_audit.md}.

\section{Next Route}

The logical next step is one bounded Refuter stress test of the audited
source-extension candidate. That stress test should ask whether
\(\XiR\) remains a coherent source-extension candidate under
extension-collapse pressure, tag-removal pressure, independence from
target calibration, and no-target-import constraints. It must preserve
all canonical ontology, source-extension adoption, benchmark, Gate
Chair, candidate reconstruction, completed-derivation, future
source-extension-impossibility, and global rejection blocks.

\begin{thebibliography}{9}

\bibitem{burden-map}
The AEther-Flow Research Project. (2026, June 17). \emph{GR derivation
burden map} [Internal control note].

\bibitem{handoff0099}
The AEther-Flow Research Project. (2026, June 18). \emph{Handoff 0099}
[Internal research-control handoff].

\bibitem{source-extension}
The AEther-Flow Research Project. (2026, June 18). \emph{Resp\_lc
source-extension candidate} [Internal research-control artifact].

\end{thebibliography}

\end{document}
